%% =====================================================================
%%  T-23-07  The Wrath of the Liar
%%  -------------------------------------------------------------------
%%  One idea per file. Core is mandatory. Slots in canonical order.
%%  Max 700 words. No layout commands. No filler. New source material
%%  for this entry goes in sources/B7/T-23-07.tex -- never by
%%  rewriting this file (docs/RULES.md R0.1).
%%  =====================================================================
%% @kind: topic
%% @id: T-23-07
%% @title: The Wrath of the Liar
%% @subtitle: Indignation is a move --- the cornered attack is itself a telling sign
%% @chapter: c23
%% @order: 70
%% @status: stable
%% @sources: B7
%% @updated: 2026-09-19

\topic{T-23-07}{The Wrath of the Liar}{Indignation is a move --- the cornered attack is itself a telling sign}

\begin{core}
B7's uncomfortable finding from watching cornered subjects, up to and including Enron's CEO on the stand: when the facts close in, the liar's most available resource is the attack --- on the questioner, the process, the audacity of the question. For the reader of behaviour, the explosion is data. For its target, the counterattack is also a pressure play: an issue with one correct answer becomes a fight about tone, where the question can be buried.
\end{core}
\begin{mechanism}
Attack mode enters the catalogue on the verbal side: backed toward facts that will not cooperate, the subject displaces energy from the issue to the meta-level --- the question was unfair, the process a witch hunt, the questioner has no right, the real victim is me. Around it: procedural complaints (objection to how replaces answer to what), sudden politeness or concern for the questioner, the simple question grown suddenly too complicated to understand. The move's power comes from its double function. As detection, indignation inside the window counts toward the cluster --- a response to the stimulus, not an answer to it. As pressure it inverts the burden: now the questioner is on trial, and defending one's manners feels more urgent than the issue; most people abandon the question there --- which trains the operator to lead with outrage next time. The defence is structural: refuse the reframe, log the attack as an answer, and re-ask. The question you were asking did not stop being legitimate because its asking offended someone.
\end{mechanism}
\begin{conditions}
Strongest where the subject's facts are genuinely bad --- the attack is proportional to the corner, which is why it spikes at the moment of the direct question. It weakens against the merely hot-tempered: some people flare at any question from anyone --- which is what the baseline is for. It also misfires where heated denial is normal register; the window's timing and the cluster's size must do the work, not volume.
\end{conditions}
\begin{application}
  \item When the counterattack lands, name its function once, internally: \emph{this is a response, not an answer}. Score it in the window.
  \item Do not defend your tone in real time. Every syllable of self-defence is a syllable not spent on the question, and both of you know the issue is dying.
  \item Re-ask, verbatim if you can, in the flattest register you own. The second asking measures the first attack.
  \item Separate the flare from the pattern: anyone can flare once. Flares timed to specific questions, absent elsewhere, are anchored to those questions.
  \item If you are the one flaring at a fair question, notice what that says. The catalogue reads its user too.
\end{application}
\begin{feedback}
  \item Working: the counterattack stopped ending conversations. You logged it and re-asked, and the second answer was more careful than the first.
  \item Working: your own flares now come with a half-second of awareness --- what question am I avoiding?
  \item Failing: you are fighting about the fight. The reframe won; the original question is dead and you are defending your politeness.
  \item Failing: every loud answer reads as guilt. You skipped the baseline; volume is only a signal against a person's own norm.
\end{feedback}
\begin{failure}
The reading's abuse is the totalitarian shortcut: define indignation as proof of guilt and no innocent person can survive an accusation --- innocence, being outraged at slander, becomes evidence. B7's own limits hold here: indignation is one indicator, worth one point in a cluster, inside the window, against a sampled baseline. A bad-faith questioner can also weaponise it: offend deliberately, then log the flare as \emph{the cluster}. The guard is the stimulus: the attack counts when it answers the question, not when it answers a provocation.
\end{failure}
\begin{countermeasures}
  \item In any dispute, watch which side answers questions and which side prosecutes the asking. The ratio is more informative than the rhetoric.
  \item Hold your own questions in a tone that gives nothing away --- the flatter the stimulus, the cleaner the window reads (T-23-05).
  \item When you feel the witch-hunt framing rising in yourself --- \emph{they are all against me} --- check whether it is answering a question you have been avoiding.
  \item Teach juniors the re-ask: it converts the oldest pressure play in interrogation into a second, better answer.
\end{countermeasures}

\sources{B7}
\seesources
\seealso{T-23-05, T-23-06, T-21-06}
